\section{Conclusion}
\label{sec:conclusion}

We presented a \emph{diagnostic and methodological} contribution for differentially private prompt optimization. On the diagnostic side, a 30-run analysis of DP-OPT on GSM8K identifies two structural failure modes of token-level DP construction: \textit{template damage}---the search trajectory deletes critical placeholders such as \texttt{\{question\}}---and \textit{noise-driven suboptimal selection}, together accounting for a 63\% failure rate. On the methodological side, we introduced \textbf{DP-ES}, a structurally cleaner evolution-based alternative that maintains a population of full prompts, mutates them without exposing private records, and restricts privacy expenditure to sampled-Gaussian scoring; selection over privatized scores is post-processing.

Across four tasks---GSM8K (88.1\%), MedQA (99.7\%), BANKING77 (73.5\%), and Alpaca (86.8\%, LLM-judge)---DP-ES matches or exceeds DP-OPT under a conservative $\varepsilon\leq1$ guarantee, with $\sim$9$\times$ lower variance and zero operational failures across 30 GSM8K runs, 2.5$\times$ faster wall-clock runtime, and 3.3$\times$ fewer logged private-data call groups. Deterministic and Gumbel-smoothed selection have comparable utility and identical privacy guarantees because both post-process privatized scores. Noise-distribution checks and a 200-profile exact-match memorization stress test provide complementary diagnostics alongside the formal RDP accounting.

\noindent\textbf{Scope.} The strongest evidence we provide is for \emph{optimization robustness under DP noise on benchmarks where prompt structure matters}, with the largest gains on reasoning-heavy tasks where token-level construction is most brittle. End-to-end validation on genuinely sensitive, non-saturated deployment data---which would require benchmarks combining record-level sensitivity with non-trivial semantic reasoning---remains future work and is the most direct extension of this contribution.
